%%%%%%%%%%%%%%%%%%%%%%% paper-emse.tex %%%%%%%%%%%%%%%%%%%%%%%%%
% TASCO Paper B 鈥?EMSE submission version (Full R1 content, EMSE format)
% Content, data, results: identical to full-paper-r1 (SUBMISSION_PREFLIGHT_PASS).
% Format adaptation only: Springer svjour3 template + spbasic (author-year) +
% abstract 150-250 words + keywords + declarations. Authors: TO BE PROVIDED.
\RequirePackage{fix-cm}
\documentclass[twocolumn]{svjour3}       % twocolumn (Springer journal format; EMSE accepts this for review)

\smartqed
\setlength{\emergencystretch}{1.5em}
\hyphenation{Large-View ex-ter-nal-i-za-tion method-ol-o-gy}
\usepackage{graphicx}
\usepackage{natbib}
\setcitestyle{round,authoryear}
\usepackage{booktabs}
\usepackage{amsmath}
\usepackage[table]{xcolor}
\usepackage[colorlinks=true,linkcolor=blue!50!black,citecolor=blue!50!black,urlcolor=blue!50!black]{hyperref}
\usepackage{balance}

\journalname{Empirical Software Engineering}

\begin{document}

\title{When Compression Backfires: An Online Dose--Response Audit of Three
Lossless Context-Compression Mechanisms for Coding Agents%
\thanks{Autonomous experimental campaign (TASCO-R4/R5/R6), human-supervised. All
artifacts, pre-registrations, and full trajectories are hash-frozen and listed in the
artifact index. This version uses only frozen evidence and adds no new online
experiment. Author information: to be provided by the research lead (see title page).}
}

\author{[AUTHORS TO BE PROVIDED BY THE RESEARCH LEAD]}
\institute{[AFFILIATION AND CORRESPONDING-AUTHOR EMAIL TO BE PROVIDED BY THE RESEARCH
LEAD. If no institutional affiliation exists, EMSE rules apply: city and country of
residence only.]}

\date{Received: date / Accepted: date}

\maketitle

\begin{abstract}
Lossless compression of tool outputs is widely assumed to reduce the token cost of LLM
coding agents. Whether removed bytes actually reduce what the model consumes --- the
only outcome that matters for cost --- is rarely tested online. We present a paired,
pre-registered online audit of three lossless compression mechanisms --- repeated-delivery
folding, eligibility-gated folding, and side-store externalization with on-demand
recovery --- across 168 disjoint tasks and 404 real agent runs of one stack (opencode;
GLM-5.3-Flash). All three mechanisms preserve task correctness, yet none yields net
input-token savings: folding 561K characters of repeated output \emph{increased}
non-cached input tokens by 7.1\%; safety gating eliminates the harm only by eliminating
the mechanism (2 fold events in 30 pairs; an offline grid over 64 eligibility
configurations shows no configuration combines $\geq$2\% safe exposure with $\geq$1/3
case coverage); and side-store externalization \emph{increases} input tokens by 23.0\%
(bootstrap 95\% CI $[-40.8, -7.6]$) because each on-demand recovery re-sends the session
context and amplifies interaction turns (+4.7 tool calls, +50.5K flow tokens per case).
We dissociate compression opportunity, triggerable exposure, and realized benefit;
attribute the inversion to compensation, recovery demand, and turn-amplified cache
traffic; and show that the safe compressible surface ($\approx$1.7\% of context) sits
below the task-level noise floor ($\hat\sigma_w \approx 0.155$, CI $[0.093, 0.197]$). A
replay-verified diagnostic audit further separates byte-level losslessness (intact) from
immediate in-context diagnostic visibility (incomplete in the deployed mechanism). We
release all protocols, 424 trajectories, the census toolkit, and the audit chain.
\keywords{coding agents \and context compression \and empirical study \and token
economics \and negative results}
\end{abstract}

\section{Introduction}
\label{intro}
Coding agents accumulate context rapidly: every tool call returns bytes that enter the
conversation and are re-sent to the model on every subsequent turn. Because trajectories
are dominated by large tool outputs --- file views, test runs, build logs, search results
--- lossless compression of this traffic is an attractive lever on token cost, and a
growing body of work compacts, summarizes, or externalizes agent context
\cite{agentdiet2026,attncompress2026,acon2026,squeez2026,taco2026}. Between an offline
compression ratio and an online bill sits a chain of confounds: the model may re-read
what was hidden (compensation), the eligibility rules that make compression safe may
shrink the compressible surface to nothing (exposure collapse), and any recovery path
re-introduces content \emph{plus} an extra interaction turn (turn amplification).

This paper reports a complete, pre-registered evidence chain that traverses this gap on
a real agent stack. We begin from an opportunity census over 2.76 billion observation
characters (20{,}000 trajectories) showing that a reversible folding candidate (RDS)
could address up to 95.96\% of opportunity mass --- the largest compressible surface we
could construct \cite{paperA2026}. We then measure, in paired online experiments, what
actually happens when that candidate --- and two successor mechanisms designed to repair
its failure modes --- are deployed on 168 disjoint tasks.

\textbf{Contributions.}
\textbf{C1: Opportunity $\neq$ triggerable exposure $\neq$ realized benefit.} A
mechanism-by-mechanism dissociation on real workloads: offline opportunity mass (561K--2.8M
characters) is far larger than what the safety contract actually lets the mechanisms
trigger online (2 fold events under gating; 280 stored views), and what triggers then
\emph{loses} 7.1--23.0\% net input tokens while preserving task correctness
(Section~\ref{sec:rq1}), with per-family structure and outliers reported in full
(Section~\ref{sec:family}).
\textbf{C2: An empirical account of benefit and its cost compensation in real agent
execution.} Three measured compensation channels invert the character savings ---
reacquisition and re-exploration (more turns under folding), recovery demand (43.4\% of
externalized views are recovered, with no pre-observable separating feature), and
turn-amplified cache traffic (+50.5K flow tokens per case; the \emph{turn}, not the
re-injected bytes, carries the cost)
(Sections~\ref{sec:rq3} and~\ref{sec:g4}).
\textbf{C3: An auditable evaluation methodology, and a diagnostic-evidence risk result.}
A pre-registered, claim-audited evaluation protocol for compression mechanisms --- paired
arms, frozen oracles, dual input calibers, an offline exposure/eligibility grid, a
recovery screen, and a byte-exact replay audit --- together with two substantive results:
the safe compressible surface is $\approx$1.7\% of context against a task-level noise
floor of $\hat\sigma_w \approx 0.155$ (CI $[0.093, 0.197]$), and keyword-caliber
diagnostic audits can pass thousands of views while semantic diagnostic lines are lost
--- byte-level losslessness and immediate diagnostic visibility are distinct properties
(Sections~\ref{sec:rq2}, \ref{sec:g4}, and~\ref{sec:g5}).
A companion measurement study characterizes the opportunity surface itself
(taxonomy, cross-corpus replication, capability coverage) \cite{paperA2026}; this
paper is its online counterpart and can be read independently.

\section{Background and Related Work}
\label{sec:related}

\textbf{Agent context economics.} LLM agents re-send their conversation on every turn,
so a tool output delivered at turn $k$ costs its tokens once at delivery and once per
subsequent turn (amortized by prompt caching). Context-compression proposals range from
summarization and truncation to external memory stores; most are evaluated by offline
compression ratios or downstream task scores. The \emph{token economics} of the
mechanism itself --- what the deployed compressor does to actual billed input --- is
rarely the dependent variable.

\textbf{Trajectory-compression mechanisms.} AgentDiet manually classifies waste in 100
Trae-Agent trajectories and reduces trajectory cost \cite{agentdiet2026} --- the closest
prior system study. AttnCompress compresses trajectories under attention guidance
\cite{attncompress2026}; Paritok-4B trains an intent-conditioned extractive compressor
over 67K OpenHands trajectories \cite{paritok4b2026}; SWE-Pruner prunes task-relevant
code context \cite{swepruner2026}; TACO self-evolves terminal-observation compression
rules \cite{taco2026}; Squeez learns minimal verbatim evidence retention for tool outputs
\cite{squeez2026}; CORVUS targets file-state synchronization to prevent stale rereads
\cite{corvus2026}; ACON optimizes context compression for long-horizon agents
\cite{acon2026}; and CoACT trains an action-preserving compressor whose compressed
observations must induce the same next action as the raw ones \cite{coact2026}. ReSum
and AgentFold apply periodic context summarization to web agents
\cite{resume2025,agentfold2025}. ICLR-style reasoning-block elision and agentic context
management systems extend the family \cite{iclr_compression2026,acm_acm2026,hasmem2026},
and TRACE studies verifier-guided compaction \cite{trace2026}. All of these answer
\emph{how to compress}; none reports paired online token economics of a lossless
mechanism across its full deployment cycle, which is our object.

\textbf{Cost and fidelity of compression.} Closest to our measurement, an empirical cost
attribution decomposes gateway-level compression costs \cite{costattribution2026}; a
fidelity study shows LLM summaries can distort decisions \cite{summariesdistort2026},
motivating preservation gates. Neither measures end-to-end paired input-token outcomes of
a frozen lossless mechanism, and neither connects behavioral compensation to the loss.

\textbf{Evaluation methodology for agents.} A 60K-trajectory stochasticity study recommends
multi-run evaluation with power analysis for sparse-activation mechanisms
\cite{randomness2026}, corroborating our noise-floor result; process-centric and failure
anatomy studies characterize behavior but not compression economics
\cite{processcentric2025,failureprocess2026}; harness-design ablations compare five
context-management strategies \cite{harnessdesign2026}; thought-action-result analyses
\cite{understandingseagents2025} and trajectory fingerprinting
\cite{trajectoriesprograms2026} complete the behavioral picture; context-rot effects
bound long-context skills \cite{contextrot2026}. Benchmarks and platforms (SWE-bench,
SWE-agent, OpenHands) evaluate task correctness \cite{swebench2024,sweagent2024,openhands2024};
we use the same stack style but measure \emph{cost}, with frozen oracles.

\textbf{Execution-effect and policy studies.} Closest to our evaluation question, three
recent studies measure what compression does to the \emph{executing} agent rather than to
single prompts. A systematic policy study varies what, when, and how much to compress
across three open-weight models on SWE-bench Verified and Terminal-Bench ($\approx$35K
runs) and finds that fewer tokens need not mean faster or cheaper execution ---
policies using roughly one-third as many tokens can take 20--80\% longer than the
uncompressed agent, with strong model and workload dependence
\cite{beyondtokensavings2026}. A controlled reacquisition-cost protocol shows that
compression can leave task completion statistically unchanged while retrieval calls
surge (21.0 $\rightarrow$ 63.9 calls, $p{=}.002$, with completion 80\%$\rightarrow$85\%),
and that the signature is environment-dependent \cite{compressioncost2026}. And at an
equal token budget, selecting whole files verbatim beats compressors and retrievers on
SWE-rebench issue resolution --- partly because compressors alter the text that patches
must reproduce \cite{skipcompressor2026}. A public benchmark (not peer-reviewed) holds a
fixed agent, model, and the official SWE-bench grader constant and ranks \emph{shipped}
compression products by cache-aware cost per solved task, reporting layers on both sides
of the no-compression baseline \cite{daseinbench2026} --- its cache-aware accounting
reinforces the point that visible-token reductions do not translate into dollar savings
when the cached prefix churns. These studies evaluate harness-level
compression \emph{policies} or shipped products; none freezes a specific lossless,
recoverable mechanism
with byte-exact recovery, pairs it against a baseline under frozen oracles, and audits
what the deployed view actually retained.

\textbf{Position.} We do not claim to be the first to evaluate compression end-to-end;
the studies above already measure execution effects at increasing breadth. Our
contribution is complementary and specific: a \emph{pre-registered, paired,
task-oracled} audit of three \emph{frozen, lossless, recoverable} mechanisms on real
repository tasks, with the failure attributed mechanism by mechanism (compensation,
exposure collapse, recovery re-injection), the safe compressible surface measured
against the task-level noise floor, and a byte-exact replay audit separating
byte-level losslessness from immediate diagnostic visibility. Positive-savings
mechanisms (AgentDiet, CoACT, ACON) remove or semantically rewrite content without
byte-exact recovery obligations; our results bound the complementary regime ---
deterministic, recoverable externalization --- where the measured outcome is a robust
loss.

\begin{figure*}[t]
\centering
\includegraphics[width=0.96\textwidth]{fig1-final.pdf}
\caption{Study overview and headline outcomes. Three lossless-recoverable compression
mechanisms (lanes) were audited in paired online runs on one stack. Badges show each
mechanism's primary observed result on its own metric: M1, input-token consumption
increase (+7.1\%, totals caliber); M2, triggered fold exposure (2 of 30 pairs); M3, mean
input-token benefit ($-23.0\%$, case-equal-weight, 95\% CI $[-40.8, -7.6]$; Table~2).
F1--F3 are mechanism interpretations supported by the execution trajectories and are not
independently proven causal conclusions. Metric definitions, denominators, and applicable
confidence intervals are given in Section~4 and Table~2.}
\label{fig:design}
\end{figure*}

\section{Study design and mechanisms}
\label{sec:mech}
All three mechanisms wrap the same frozen primitive family and share the deployment
contract: anchor-first delivery, exact-content-only folding, fail-closed behavior on any
inconsistency, auditable markers, and byte-exact reconstruction
(Figure~\ref{fig:design}). Losslessness here means byte-exact reconstruction from the
side store or from re-delivered anchors; it does not by itself guarantee that every
diagnostically relevant line is \emph{directly visible} in context --- the two properties
are distinct (Section~\ref{sec:g5}).

\textbf{M1 --- Repeated-delivery folding (RDS, high dose).} Any contiguous run of
$\geq$2 lines exactly matching a contiguous run delivered in a \emph{single earlier}
observation is replaced by an audit marker; the first delivery is always verbatim
(anchor-first). Deployed at full coverage in R4: every tool output in every run.

\textbf{M2 --- Eligibility-gated folding (safe dose).} The same primitive with three
gates: diagnostic-bearing outputs are never folded; target/specification file reads are
never folded; read outputs fold only from the third delivery of the same path; non-read
outputs fold only if $\geq$1{,}000 characters with duplicate-line share $\geq$0.30. Gates
use only runtime-observable facts.

\textbf{M3 --- Side-store externalization (LargeView).} Read outputs $\geq$4{,}000
characters are replaced in-context by head (20) and tail (20) windows, verbatim copies
of lines matching a frozen eight-keyword diagnostic pattern
(\texttt{fail\allowbreak|error\allowbreak|assert\allowbreak|traceback\allowbreak|exception\allowbreak|fatal\allowbreak|denied\allowbreak|not found}), and a pointer
marker; the full original is written losslessly to \texttt{.r5-store/}\allowbreak\texttt{<sha16>.txt} for
recovery by an ordinary read. \emph{Post-hoc disclosure:} this keyword caliber does not
retain failure returns that carry no listed keyword, and markers print raw tool-output
line numbers while the view displays content line numbers; both defects are bounded,
deterministic, and characterized in Section~\ref{sec:g5}.

Each mechanism is deterministic, requires no model calls, and was frozen (code hash,
protocol, oracle set) before its online evaluation.

\section{Methodology}
\label{sec:method}

\textbf{Frame and stack.} Four experiments over 168 disjoint tasks drawn from two real
JavaScript repositories, executed by the opencode harness \cite{opencode2026} driving
GLM-5.3-Flash (model identity verified per run from harness logs). Tasks: code review,
module analysis, test-fix/TDD, and cross-specification analysis --- selected by frozen
seeded rules from unused material; no task was chosen or dropped after observing
results. Case-level disjointness across the four experiments is verified programmatically
(zero pairwise case-ID intersections); within M3, the two repetitions per arm are
repeated measures on the same cases, aggregated to arm means \emph{before} pairing (60
runs are not 60 independent tasks). The LLM under study serves as the execution engine
of the audited stack; it is not an author of this work, and its use is documented here
and per run in the harness logs.

\textbf{Design.} Paired per case: every case executes once with compression disabled
(Baseline, B) and once enabled (Treatment, T), in seed-balanced AB/BA order, with a fresh
copy of the pinned repository per run. M3 adds a second repetition per arm (2 reps
$\times$ 2 arms). A baseline-only noise pilot (Stage-0) measured within-task variance on
8 tasks $\times$ 3 repetitions.

\textbf{Calibers (pre-registered).} Table~\ref{tab:calibers}. \emph{Input tokens}
(non-cached, api-reported; primary), \emph{cache-read tokens}, \emph{output tokens},
\emph{flow tokens} (sum of the former three plus reasoning), \emph{wall time}, \emph{tool
calls}, \emph{oracle success} (frozen test commands or deterministic text-token oracles),
\emph{fold/store events}, and \emph{recovery reads}. Character counts are opportunity
proxies only and are never added to token quantities.

\begin{table}[t]
\centering\small
\caption{Measurement calibers and their cautions.}
\label{tab:calibers}
\begin{tabular}{@{}p{1.9cm}p{3.6cm}p{2.1cm}@{}}
\toprule
Caliber & Definition & Caution \\
\midrule
Chars saved & obs chars removed by fold/store & proxy; not tokens \\
Input & non-cached input tokens (api-reported) & primary metric \\
cache\_read & cached re-send tokens & amplifies with turns \\
Flow & in+out+\allowbreak reasoning+\allowbreak cache & session-length sensitive \\
Success & oracle passed & B/T same oracle \\
\bottomrule
\end{tabular}
\end{table}

\textbf{Statistics.} Case-level pairing. We report two pre-declared input aggregates and
never mix them: a \emph{totals-based} delta, $(\sum T - \sum B)/\sum B$, which weights
cases by their size (the historical headline caliber for M1), and a
\emph{case-equal-weight} mean of per-case relative deltas with case-level bootstrap 95\%
CIs (10{,}000 resamples; the headline caliber for M3). Robustness via medians,
win/loss/tie counts, and leave-one-out sensitivity. Families with fewer than 10 cases are
marked \emph{exploratory}. Uncertainty via case-level bootstrap; the Stage-0 noise floor
is $\hat\sigma_w = 0.155$ (task-level bootstrap 95\% CI $[0.093, 0.197]$), and the
expansion directly measured $\hat\sigma_w(T) = 0.258$ vs.\ $\hat\sigma_w(B) = 0.244$
(consistent with equality at this sample size --- not a proof). A pre-registered 15-claim
audit with evidence grades (Section~\ref{sec:claims}) governs what each result may be
used to claim.

\textbf{Offline instruments.} Three deterministic, zero-model-call instruments close the
design: (i) an eligibility \emph{grid} that replays all 100 baseline trajectories under 64
gate configurations; (ii) a \emph{recovery screen} that censuses every externalized view
for recoveries, distances, and per-view character nets; (iii) a \emph{replay audit} that
reconstructs every compressed view byte-for-byte from frozen artifacts (used by the
Section~\ref{sec:g5} case study).

\section{Results}

\subsection{RQ1: Does lossless compression reduce real input tokens? No.}
\label{sec:rq1}

Table~\ref{tab:unified} unifies the three mechanisms. Figure~\ref{fig:unified} shows the
primary metric; Figure~\ref{fig:dists} the per-case distributions.

\begin{table*}[t]
\centering\small
\setlength{\tabcolsep}{4pt}
\caption{Unified mechanism table (columns = mechanisms; calibers never mixed).
``Input $\Delta$ (totals)'' weights cases by size; ``benefit (case mean)'' is the
case-equal-weight mean of per-case relative deltas with bootstrap 95\% CI. Benefit
values are color-coded: \textcolor[rgb]{0.55,0.09,0.18}{red} = input loss,
\textcolor[rgb]{0.10,0.42,0.22}{green} = saving.}
\label{tab:unified}
\begin{tabular}{@{}p{3.55cm}p{4.35cm}p{4.35cm}p{4.35cm}@{}}
\toprule
Metric & \textcolor[rgb]{0,0.45,0.70}{\textbf{M1 folding}} & \textcolor[rgb]{0.80,0.47,0.04}{\textbf{M2 gated}} & \textcolor[rgb]{0.0,0.47,0.42}{\textbf{M3 side-store}} \\
 & \emph{100 pairs, 1 rep/arm} & \emph{30 pairs, 1 rep/arm} & \emph{30 cases, 2 reps/arm} \\
\midrule
Online exposure & 542 folds; 561K chars; 64/100 cases & 2 fold events; 2/30 cases & 280 views; 2.81M chars; 123 rec. \\
Input $\Delta$ (totals) & $+7.1\%$ & $-2.9\%$ & $+12.4\%$ \\
\rowcolor[rgb]{0.96,0.96,0.96} Benefit, case mean [95\% CI] & \textcolor[rgb]{0.55,0.09,0.18}{$-41.8\%$} [$-66.7$, $-20.2$] & \textcolor[rgb]{0.55,0.09,0.18}{$-21.3\%$} [$-50.3$, $+2.6$]$^{a}$ & \textcolor[rgb]{0.55,0.09,0.18}{$-23.0\%$} [$-40.8$, $-7.6$] \\
\rowcolor[rgb]{0.96,0.96,0.96} Median benefit & $-4.0\%$ & $+4.9\%$ & $-15.3\%$ \\
Cache B$\rightarrow$T & 10.27$\rightarrow$10.74M (+4.6\%) & 5.36$\rightarrow$4.42M ($-17.6\%$) & 2.20$\rightarrow$3.58M (+63.2\%) \\
Turns / tool calls & +50 harness steps (total) & tool-call ratio 1.02 & $+4.7$ calls; $+2.4$ turns \\
Task B$\rightarrow$T (oracle) & 81$\rightarrow$80 & 24$\rightarrow$24 & 18$\rightarrow$20 \\
Notable case & C-ENT04 (T input $8.7\times$ B) & R-XS-AUTH-REP ($-273\%$) & X-HI-M01 ($+23.9\%$, 0 rec.) \\
\bottomrule
\multicolumn{4}{@{}p{\textwidth}@{}}{\footnotesize $^{a}$M2's case-mean is dominated by one exploratory outlier; its median is $+4.9\%$, and its paired flow-delta median is $-0.25\%$ --- the mechanism is near-inert, as the grid predicts. M1's case-mean is outlier-heavy too; medians give the robust picture --- still at or below parity. Success rows are observed oracle parity, not a powered non-inferiority test. M3 wall $+48\%$.}
\end{tabular}
\end{table*}

\textbf{M1 (100 pairs):} 542 fold events removed 561{,}454 characters, yet treatment arms
consumed \textbf{7.1\% more} non-cached input than baseline (totals caliber; win/loss on
flow 35/62). Task success showed no observed degradation (81\% $\rightarrow$ 80\%); this
parity is an observation, not a powered non-inferiority test.

\textbf{M2 (30 pairs):} gates eliminated the harm (median paired flow delta $-0.25\%$,
tool-call ratio 1.02, success 24/24 both arms) --- but also eliminated the mechanism: only
2 fold events fired across 30 cases.

\textbf{M3 (30 cases $\times$ 2 reps):} 280 stored views (2.81M characters externalized),
123 recovery reads; treatment arms consumed \textbf{23.0\% more} non-cached input
(case-equal-weight; bootstrap 95\% CI $[-40.8\%, -7.6\%]$; median $-15.3\%$; win/loss
8/22). Success 18/30 $\rightarrow$ 20/30. Across 168 disjoint tasks and three designs, no
configuration produced positive net input savings, and the two mechanisms with measurable
exposure produced losses whose bootstrap 95\% CI excludes zero
(Figure~\ref{fig:unified}).

\begin{figure}[t]
\centering
\includegraphics[width=\columnwidth]{f1_unified_benefit.pdf}
\caption{Primary metric by mechanism: case-equal-weight mean input benefit with bootstrap
95\% CI (bars), and totals-based input delta (diamonds; positive value = treatment used
more input). Zero line = parity. All three mechanisms are at or below parity.}
\label{fig:unified}
\end{figure}

\begin{figure*}[t]
\centering
\includegraphics[width=0.92\textwidth]{f2_paired_deltas.pdf}
\caption{Per-case paired input benefit distributions (points jittered; box = quartiles;
zero = parity). Distributions are wide relative to their means --- the noise floor of
Section~\ref{sec:rq2}.}
\label{fig:dists}
\end{figure*}

\subsection{Family structure and robustness}
\label{sec:family}

\begin{table}[t]
\centering\scriptsize
\caption{Per-family input benefit (case-equal-weight mean, 95\% CI, win/loss) by
mechanism. $n<10$ = exploratory. Families are separate tasks, not replications.}
\label{tab:family}
\setlength{\tabcolsep}{3.2pt}
\begin{tabular}{@{}llrrl@{}}
\toprule
Family & Mech & $n$ & Benefit [CI] (\%) & W/L \\
\midrule
file-review & M1 & 48 & \textcolor[rgb]{0.55,0.09,0.18}{$-59.9$} [$-95.3,-29.0$] & 17/31 \\
            & M2 & 8$^{e}$ & \textcolor[rgb]{0.55,0.09,0.18}{$-7.7$} [$-43.4,+17.2$] & 4/4 \\
            & M3 & 18 & \textcolor[rgb]{0.55,0.09,0.18}{$-32.9$} &  \\
module-analysis & M1 & 24 & \textcolor[rgb]{0.55,0.09,0.18}{$-6.9$} [$-22.2,+6.1$] & 12/12 \\
            & M2 & 8$^{e}$ & \textcolor[rgb]{0.55,0.09,0.18}{$-28.5$} [$-96.5,+16.9$] & 5/3 \\
            & M3 & 8$^{e}$ & \textcolor[rgb]{0.55,0.09,0.18}{$-9.9$} &  \\
mutation-testfix & M1 & 15 & \textcolor[rgb]{0.55,0.09,0.18}{$-11.1$} [$-38.8,+11.0$] & 6/9 \\
            & M2 & 8$^{e}$ & \textcolor[rgb]{0.55,0.09,0.18}{$-3.6$} [$-26.5,+19.6$] & 5/3 \\
            & M3 & 2$^{e}$ & \textcolor[rgb]{0.55,0.09,0.18}{$-39.3$} &  \\
spec/contract & M1 & 10 & \textcolor[rgb]{0.10,0.42,0.22}{$+7.1$} [$-21.2,+33.0$] & 7/3 \\
cross-spec & M2 & 3$^{e}$ & \textcolor[rgb]{0.55,0.09,0.18}{$-87.3$} [$-272.9,+11.7$] & 1/2 \\
           & M3 & 2$^{e}$ & \textcolor[rgb]{0.10,0.42,0.22}{$+30.3$} &  \\
enterprise & M1 & 3$^{e}$ & \textcolor[rgb]{0.55,0.09,0.18}{$-346.7$} [$-774.7,-18.4$] & 0/3 \\
\bottomrule
\end{tabular}
\end{table}

Table~\ref{tab:family} breaks the primary metric down by task family
(Figure~\ref{fig:family}). Cells with $n<10$ are exploratory signals, not formal tests;
three observations. (i) \emph{No family escapes the negative
sign where exposure is measurable}: file-review --- the family with the most folding
opportunity --- is the worst under M1 ($-59.9\%$ mean, CI excluding zero), exactly where
compensation predicts the largest harm. (ii) \emph{Positive cells are small, exploratory,
and mechanism-consistent with near-zero exposure}: M1 spec/contract ($+7.1\%$, $n{=}10$,
CI includes zero) and M3 cross-spec ($+30.3\%$, $n{=}2$) have the least exposure. (iii)
\emph{Oracle parity hides family-level movement}: M1 spec/contract success moves
4$\rightarrow$2/10 and M2 mutation 8$\rightarrow$7 --- single-family counts too small for
inference, reported for completeness. Robustness: M3's headline survives leave-one-out
($-17.2$pp without the largest outlier, X-NO-F01; median $-15.3\%$; win/loss 8/22), and a
zero-exposure control subgroup (6 cases with no markers; mean $+10.6\%$, median $+2.4\%$)
is statistically indistinguishable from zero, excluding systematic T-arm penalties such
as plugin loading or ordering effects. M1/M2 means are outlier-sensitive
(leave-one-out ranges in the supplement); their medians ($-4.0\%$ / $+4.9\%$) are the
robust summary --- still at or below parity.

\begin{figure*}[t]
\centering
\includegraphics[width=0.96\textwidth]{f6_family_benefits.pdf}
\caption{Per-family input benefit (mean, bootstrap 95\% CI) by mechanism. Exploratory
cells ($n<10$) are marked. No family with measurable exposure escapes the negative sign.}
\label{fig:family}
\end{figure*}

\subsection{RQ2: How large is the \emph{safe} compressible surface? Below measurability.}
\label{sec:rq2}

The exposure--safety tension is not a tuning artifact. An offline grid replayed
all 100 baseline trajectories under 64 eligibility configurations (read-delivery
threshold $\times$ duplicate-share threshold $\times$ minimum size; 70/30 calibration/holdout
split). \textbf{No
configuration} satisfies the pre-registered acceptance gates of $\geq$1/3 case coverage
and $\geq$2\% safe exposure (Figure~\ref{fig:grid}): the Pareto frontier tops out at
\textbf{1.66\%} exposure (the most permissive safe configuration, which is exactly the M1 harm regime minus protected content), and the configuration
matching M2's gates replays to \textbf{zero} fold events on the same trajectories while
its protection gates block 245 target/specification/diagnostic folds --- explaining M2's
near-inertness directly. The safe surface is structurally below the task-level noise
floor: at $\hat\sigma_w \approx 0.155$, no replication count $\leq$10/arm reaches 0.80
detection probability for effects of $+3$--5\% \cite{randomness2026}.

\begin{figure}[t]
\centering
\includegraphics[width=\columnwidth]{f4_grid_pareto.pdf}
\caption{Offline eligibility grid (64 configurations, holdout replay): safe exposure vs.\
case coverage. Dashed/dotted lines: pre-registered acceptance gates ($\geq$2\% exposure,
$\geq$1/3 coverage). No configuration passes; the frontier tops out at 1.66\%.}
\label{fig:grid}
\end{figure}

\textbf{Cross-corpus exposure.} A second-workload census (15{,}000 sessions from an
independent corpus) confirms that the \emph{exposure} pattern generalizes: side-store
eligible views trigger 52{,}719 times, saving 366M characters --- 29.5\% of all
observation characters there. Exposure is not the bottleneck on that corpus; its
\emph{recovery economics} remain untested online. Both instruments are character-level
opportunity proxies on recorded data --- they cannot establish online token benefit,
which is precisely the lesson of this paper.

\subsection{RQ3: Can eligibility features predict recovery need in advance? No.}
\label{sec:rq3}

Across 311 externalized views (70 treatment runs), 135 (43.4\%) were recovered at least
once (137 recovery events; median distance 4 tool events after the fold;
Figure~\ref{fig:cdf}). Recovery rates by size bucket, file extension, task family, and
exposure stratum all fall in the 38--52\% band with no separating feature, in both the
design split and a 30\% holdout (43.4\% vs.\ 49.3\%). Per-view character nets remain
positive even including recoveries (design $+4{,}241$, holdout $+4{,}589$ chars/view) ---
yet the same cases lose 23\% input tokens. The only strongly positive exposed case
(X-HI-M01: $+23.9\%$) had \textbf{zero} recoveries --- the winning condition (``the model
never needs the content back'') is observable only post hoc. A conservative caveat: 24
same-file rereads occurred in baseline arms without any side store, so ``no recovery''
does not mean ``no need''; needs can resurface as re-derivation, not reads.

\begin{figure}[t]
\centering
\includegraphics[width=\columnwidth]{f5_recovery_cdf.pdf}
\caption{CDF of the event distance from a fold to its first recovery (recovered views,
$n{=}135$; median 4). Recoveries come early, re-pricing the session while its context is
largest.}
\label{fig:cdf}
\end{figure}

\subsection{Why character savings become token losses: turn amplification}
\label{sec:g4}

The central measurement lesson is that \emph{character accounting and API token
accounting are separate calibers, and their signs diverge}. We therefore report them
separately and never add them.

\textbf{Caliber A --- API tokens (primary, api-reported).} Per case, M3 treatment arms
consumed $+2{,}577$ more non-cached input tokens, $+46{,}265$ more cache-read tokens
($+50.5$K flow tokens), issued $+4.7$ tool calls and $+2.4$ model turns, and took 48\%
longer. The dominant observed cost is cache-read traffic, which scales with session
length and turn count: $\Delta$cache tracks $\Delta$flow almost perfectly
($r{=}0.997$), and $\Delta$cache correlates with $\Delta$input at $r{=}0.33$.

\textbf{Caliber B --- mechanism characters (measured, arm means).} The primitive saved
$S = 650{,}693$ characters, re-injected $R = 158{,}880$ through recovery reads, and added
$M = 40{,}600$ marker characters: a net \emph{character} saving of $+451{,}214$. The
per-case decomposition leaves a residual $X = 528{,}537$ characters that absorbs the
unit conversion and treatment-arm behavioral divergence (output $+13\%$, wall $+48\%$);
we do not interpret $X$ causally. \emph{Caliber note:} recovery re-injection also exists
in a per-run caliber (352{,}804 characters over all 70 treatment runs including smoke
runs, or 317{,}759 over the 60 expansion runs); the arm-mean caliber above is the one
compatible with per-case pairing. Both are frozen and reconciled in the artifacts.

\textbf{The divergence.} Why does a $+451$K-character saving coexist with worse
API-token outcomes? Because the binding cost is not the returned bytes but the
\emph{turn}: each recovery is mechanically one additional tool call whose re-sent context
is billed largely as cache reads. Consistently, the \emph{size} of what comes back
barely predicts the loss (re-injected chars vs.\ per-case input delta: $r{=}0.03$),
while \emph{how much extra traffic the session generates} does
($r(\Delta\text{cache},\Delta\text{input}) = 0.33$). Treatment trajectories also diverge
behaviorally (exploration, output $+13\%$), which we measure but cannot causally
decompose without counterfactual arms; we therefore claim only the association, in both
directions. Per-view character accounting is structurally blind to this effect; any
compressor with a recovery path is exposed to it (Figure~\ref{fig:amp}).

\begin{figure}[t]
\centering
\includegraphics[width=\columnwidth]{f3_turn_amplification.pdf}
\caption{Turn amplification: $\Delta$ tool calls vs.\ $\Delta$ flow tokens (K) per case;
bubble size = recoveries. More recoveries $\Rightarrow$ more turns $\Rightarrow$ more
re-sent context.}
\label{fig:amp}
\end{figure}

\subsection{A bounded diagnostic-visibility case study: the post-hoc G5 audit}
\label{sec:g5}

\textbf{Motivation and method.} The original claim that M3 retains ``all diagnostic lines
verbatim'' had been verified only in a \emph{keyword caliber}: a frozen audit checked
that keyword-matching line counts are identical before and after folding across 24{,}648
replayed views. A post-hoc audit (G5) tested the \emph{semantic} caliber on 30 blind
views drawn from 3 evidence runs (10/13/7 views; 3 underlying cases). Two isolated LLM
reviewers reviewed the blind package under a frozen protocol; every verdict-driving
claim was then re-verified mechanically, so the findings below do not rest on model
votes. \emph{This is a model-assisted audit; independent human review remains pending.}

\textbf{Findings (all mechanically verified by deterministic replay).}
(1)~\emph{E1, semantic diagnostic loss from the view}: 12 unique failure-return lines
across the 3 target files (4 per file,
\texttt{return \{ ok: false,\allowbreak reason: "exhausted\_retries",\allowbreak attempts: attempt \};}),
appear in none of the 30 in-context views; they are retained
only in the byte-exact side store. The deployed keyword pattern does not match them.
The store itself is faithful: stored content equals the repository files line-for-line
(30/30), marker hashes match, and each in-context view is byte-identical to the
deterministic replay of the frozen folding function over the stored original (30/30) ---
the view is fully explained by the mechanism, and the agent session streams contain
exactly the folded views. (2)~\emph{E3, marker coordinate inconsistency}: markers
declare the omitted range in \emph{raw tool-output} line numbers
(\texttt{21..198/199}) while the surrounding view displays \emph{content} line numbers;
the 3-line read envelope shifts the two systems by a constant, so a reader reconciling
marker and view concludes that content lines 18--20 are unaccounted. In content
coordinates the declared range exactly abuts the head and tail windows: no content is
undeclared-omitted; the defect is one of
auditability, not of content loss.

\textbf{Scope and status.} The audit covers 3 evidence runs (3 cases, 30 views); it does
not estimate a population defect rate, and the 30 views are not 30 independent tasks.
The two findings are deterministic properties of the frozen plugin (one keyword-list
coverage gap; one coordinate-system choice), not random failures. The economic results
above are measured in token calibers and are unaffected; oracle correctness is
unaffected. The claim ``all diagnostic lines verbatim'' is \textbf{withdrawn} in favor
of the keyword-caliber statement in Section~\ref{sec:mech} plus this case study;
byte-level losslessness and recovery integrity were independently re-verified and stand.
A repair design (semantic diagnostic protection with structure-aware failure-return
detection and coordinate-consistent markers) is released with the artifacts; a repaired
mechanism is a \emph{new} mechanism and would require fresh online measurement.

\section{Discussion: a reusable failure model}
\label{sec:disc}

The three mechanisms fail in three distinct, now-measured ways, and the failures
generalize into a pre-deployment checklist for any context compressor:

\textbf{F1 --- Compensation (M1).} Folding hides content the model was using: harmed
families explore \emph{more} (file-review: +24\% total harness steps under M1, computed
from frozen per-case data), consistent with
agents re-acquiring hidden context. \emph{Gate: measure compensation before deployment;
folding is not free just because it is reversible.}

\textbf{F2 --- Exposure collapse (M2).} Safety gates remove exactly the folds that caused
F1 --- and with them 98\% of the exposure; the grid shows no parameter point escapes the
frontier. \emph{Gate: run the exposure census first --- if the safe surface is below a
few percent of context, the mechanism cannot produce a measurable effect.}

\textbf{F3 --- Recovery re-injection (M3).} Side-store recovery restores full content
natively plus one extra turn; recoveries come early (median 4 events), re-pricing the
session when its context is largest, and no pre-observable feature separates
low-recovery views. \emph{Gate: price recovery as turns, not bytes; require a
turn-neutrality argument for any call-back path.}

\textbf{For practitioners.} (1) Run the exposure census --- opportunity proxies are
necessary but not sufficient; (2) price recovery in turns and cache traffic, not
characters; (3) measure per-turn amplification, not per-fold savings; (4) verify
correctness with the same oracle as the deployment target; (5) verify \emph{diagnostic
visibility} against a semantic definition of ``diagnostic,'' not a keyword list --- our
keyword-caliber audit passed 24{,}648 views while missing every non-keyword failure
return (Section~\ref{sec:g5}).

\textbf{For researchers.} The exposure--measurability gap generalizes the
variance-dominance problem of agent evaluation: mechanism exposure must exceed the
within-task noise floor ($\hat\sigma_w \approx 0.155$ here) by the detection margin, or
online experiments are underpowered regardless of sample size. Our census toolkit
computes both quantities offline.

\textbf{Applicability boundary.} The three failure mechanisms were measured on
\emph{deterministic, lossless, recoverable} mechanisms on multi-turn repository tasks.
They do not extend to the positive-savings regimes reported elsewhere: removal-based
reduction \cite{agentdiet2026}, learned action-preserving compression \cite{coact2026},
and guideline-semantic compression \cite{acon2026} trade byte fidelity (and recovery
obligations) for savings, and operate in settings our exposure analysis does not cover.
Our compensation findings are nonetheless consistent with independent protocols ---
reacquisition surges with completion unchanged \cite{compressioncost2026}, and
policy-level studies find fewer tokens can mean slower, costlier execution
\cite{beyondtokensavings2026} --- suggesting the compensation channel is not specific to
our stack, while the \emph{magnitude} of every number here is.

\section{Pre-registered claim audit}
\label{sec:claims}

Fifteen claims were pre-registered with evidence grades: 7 STRONGLY\_SUPPORTED, 4
SUPPORTED\_WITH\_LIMITATIONS, 2 HYPOTHESIS\_ONLY, and 2 UNSUPPORTED (the latter two are
overreach formulations we explicitly disavow, including ``all compression strategies are
useless''). The full table with evidence pointers is in the supplement (Online Resource
1); the load-bearing
entries: C1 no net input gain (SS), C2 turn amplification dominates (SS), C3 safe surface
$\approx$1.7\% (SS), C4 no correctness regression observed (SS, caliber-limited), C7
offline census predicts exposure (SS), C8 $\hat\sigma_w = 0.155$ $[0.093,0.197]$ (SS), C9
recovery is real agent behavior (SS), C12
turn-neutral silent truncation might work \emph{--- only under an explicit
diagnostic-visibility guarantee; the deployed M3 lacked one (\S5.7)} (H), C13
semantic-fidelity compression viable \emph{--- hypothesis only; the deployed
keyword-caliber diagnostic protection did not achieve semantic diagnostic visibility
(\S5.7)} (H).

\section{Threats to validity}
\label{sec:threats}
\textbf{External validity.} One model (GLM-5.3-Flash), one harness (opencode), one
language family (JavaScript repositories), and analysis-type tasks dominate the sample;
the online evidence base is a single stack. A systematic policy study finds compression
effects are strongly model- and workload-dependent \cite{beyondtokensavings2026}, so our
numbers should be read within that dependence. The second-workload census is offline
(character proxies only), so its recovery economics are untested. Mutation-heavy
workloads showed directionally positive subgroups under M1/M3 but are small ($n{=}15$,
$n{=}8$) and were not powered.
\textbf{Statistical validity.} M3 uses 2 reps/arm aggregated to arm means (30 independent
cases); M1/M2 use 1 rep/arm. Case-level bootstrap CIs are reported for all headline
means; family cells with $n<10$ are exploratory. We report both totals-based and
case-equal-weight input calibers and do not mix them.
\textbf{Construct validity.} Oracle success is not diagnostic visibility; the two are
kept separate (Section~\ref{sec:g5}). Character counts are opportunity proxies.
\textbf{Measurement and cost.} Harness-side token accounting is verified against
api-reported usage, but billing amounts were not observable, and we report no
estimated-cost conversions --- all cost statements are API-observable metering.
A tokenizer-exact, per-recovery causal
decomposition is not observable from harness telemetry; the $S/R/M/X$ decomposition is
an accounting identity whose residual absorbs the unit conversion, and we do not
interpret it as a physical identity.
\textbf{Safety-claim scope.} The fail-closed contract is verified by code review, unit
and gate tests (17/17 primitive semantics, 9/9 gates), and a zero-incident observation
across 424 runs; no systematic fault-injection campaign was executed. Safety claims are
therefore limited to these tested invariants and observed behavior, not to arbitrary
fault hypotheses.
\textbf{Internal validity.} B/T trajectories diverge before any marker in companion
analyses, so single-run deltas conflate variance with mechanism; the zero-exposure
control subgroup and multi-experiment consistency mitigate but do not eliminate this.
\textbf{Audit and mechanism-verification scope.} The Section~\ref{sec:g5} audit is
model-assisted (two isolated LLM reviewers whose verdicts diverged substantially and
were therefore not used as evidence; all reported defects verified by deterministic
replay). Independent human review is registered but pending; the audit covers 3 evidence
runs (3 cases), and only those views are byte-replay-verified --- no population
defect-rate is inferred, and the remaining runs are covered by the mechanical preflight
(direction-inclusive containment) rather than semantic audit.

\section{Conclusion}

We built the largest safe compressible surface we could construct for a real coding
agent, deployed it online under a pre-registered paired design, and measured a robust
negative. The measurement chain --- offline census, dose--response, gating, recovery
accounting, and a replay-verified diagnostic audit --- explains why, and the explanation
generalizes across the three frozen mechanisms studied here: for lossless compression on
agentic workloads, the binding constraint is not compression ratio but the economics of
interaction turns and the unpredictability of content re-use. This negative economics is
a property of the measured v1 mechanisms; it does not by itself prove that every
conceivable repair must lose --- a diagnostic-safer externalization (semantic protection
with coordinate-consistent markers; repair design released) changes the exposure
trade-off and must be judged by the same online measurement protocol before any such
claim. We release the census toolkit, all protocols, 424 trajectories, the
diagnostic-audit evidence chain, and the repair design to make the next mechanism's
evaluation cheaper than ours was.

\section*{Declarations}

\textbf{Funding.} [TO BE PROVIDED BY THE RESEARCH LEAD --- or state ``No funding was
received for this research.'']

\textbf{Competing interests.} [TO BE PROVIDED BY THE RESEARCH LEAD --- at minimum
disclose that one of the audited compression mechanisms originates from the same
research programme as this paper, if applicable.]

\textbf{Data availability.} [TO BE FINALIZED BY THE RESEARCH LEAD. Factual basis: the
complete evidence chain --- 424 agent-run trajectories, events, protocols, pre-registrations,
analysis scripts, and the audit chain --- exists and is hash-frozen (artifact index,
1{,}889 files across R3--R5). Options: (a) public repository DOI of the frozen index;
(b) available from the corresponding author on request. Second-corpus census results are
aggregate counts over public datasets \cite{thoughtworks2026,openswetraces2026}.]

\textbf{Code availability.} [TO BE FINALIZED BY THE RESEARCH LEAD. The analysis scripts
(bootstrap, replay audit, census) are deterministic and re-runnable; packaging and
release channel per (a)/(b) above.]

\textbf{AI tools.} The studied coding-agent stack executes an LLM (GLM-5.3-Flash served
via a Volcano-Engine-compatible gateway \cite{opencode2026}); it is the \emph{object} of
measurement, not an author. LLM-based tooling supported parts of the analysis
automation; every reported number was verified by deterministic replay over frozen
artifacts, and human accountability for the final text rests with the authors.

\textbf{Scope boundary with the companion study.} Paper~A \cite{paperA2026} measures the
compression-\emph{opportunity} surface (taxonomy, replication, capability coverage) on
public corpora; this paper reports the \emph{online economics} of three frozen
mechanisms on its own paired experiments. All online numbers in this paper come from
Paper~B's own runs; Paper~A's census figures are cited only as prior characterization
and are never counted as this paper's experimental results.

\begin{appendix}
\section{Supplement}
The supplement (Online Resource 1, \texttt{ESM\_1.pdf}) contains: the full caliber
table, Stage-0 noise detail, the 64-configuration eligibility grid, robustness
detail (win/loss, leave-one-out), the R6 recovery-screening detail, the G5 audit
protocol and limited-sample notes, the repair-design summary, the C1--C15 claim table
with evidence pointers, and the number-to-evidence map.
\end{appendix}

\balance
\bibliographystyle{spbasic}
\bibliography{refs}

\end{document}
